\section{Introduction}
\label{sec:intro}

\paragraph{Motivation.}
A central goal in coding theory and in cryptography-facing applications is to obtain linear minimum distance together with
streaming encodability, i.e., the ability to compute the encoding in a single pass over the input while producing output symbols incrementally and using only small working memory \stan{Peter: check if i explained correctly}, and near-linear work.  Classical turbo-style constructions suggest that random interleaving
and recursion can produce excellent distance behavior in practice, but converting this intuition into
clean, modular, proof-driven guarantees, especially under stringent work and locality constraints, remains subtle.

\paragraph{Our viewpoint: serial concatenation as a modular distance amplifier.}
We study depth-2 serial concatenation with a random interleaver:
\[
\Csc \;:=\; \Ctwo \circ \Pi \circ \Cone \subseteq \F_2^n,
\qquad
G_{\mathrm{sc}} \;=\; \Gone \cdot \Pi \cdot \Gtwo,
\]
where \(\Cone\) is an outer code, \(\Pi\) is a uniform random interleaver on \([n]\), and \(\Ctwo\) is a rate-1 recursive inner encoder. \stan{we need to add what G is}
Our main conceptual contribution is to separate the minimum-distance proof into two independently verifiable interfaces:

\begin{itemize}
\item \textbf{Outer interface:} (i) logarithmic minimum distance \(d_{\min}(\Cone)\ge \alpha\log n\), and
(ii) a spectrum bound \(\enum_{\Cone,w}\le \poly(n)\beta^w\) on the weight regime used by the first-moment argument.
\item \textbf{Inner interface:} a slice-to-tail bound for a random interleaved input of weight \(w\),
\[
p_w(\delta)\ :=\ \Pr\big[\wt(\Ctwo(U))\le \delta n \,\big|\, \wt(U)=w\big],
\]
where \(U\) is uniform over the Hamming slice \(\{0,1\}^n\) conditioned on \(\wt(U)=w\).
\end{itemize}

Once these interfaces are established, a short first-moment calculation yields linear minimum distance for \(\Csc\) with high probability
(over \(\Pi\) and any inner randomness).

\stan{Add to prelims - spectrum bound, first moment argument}
\stan{seems like standard terms, but i had to parse the second bullet point with chatgpt}

\paragraph{Target performance characteristics.}
Our design goals are the following.
We want (i) provable linear minimum distance \(d_{\min}(\Csc)=\Omega(n)\) with high probability over the randomized construction
(interleaver and any internal randomness), (ii) streaming/online encoding with small working memory (ideally \(O(\log n)\) or \(O(1)\) state),
and (iii) near-linear work, measured by tap complexity: the number of state coordinates read/XORed per output symbol.
Concretely, we aim for \(O(1)\) taps per symbol (hence \(O(n)\) total taps) with a single random interleaver.
Our modular framework makes these tradeoffs explicit: the outer code controls the low-weight spectrum, the interleaver randomizes support,
and the inner recursion supplies contraction  and large-deviation behavior.


\paragraph{Outer instantiations.}
We instantiate the outer interface in one of two ways.
\begin{enumerate}
\item \textbf{Dense log-memory outer.}
A terminated, feedforward convolutional outer with memory \(M=\Theta(\log n)\) and dense taps within the band.
It is encoding-friendly and satisfies both the logarithmic distance and the required spectrum shape.
\item \textbf{Encodable banded LDPC outer (local expander on a cycle).}
A banded LDPC parity-check matrix \(H=[A\mid B]\) with locality window \(W=o(n)\), where \(B\) is a staircase+gap parity part
(enabling linear-time encoding) and \(A\) contributes random local edges that yield local expansion.
This outer achieves the same interface guarantees while supporting linear-work encoding.
\end{enumerate}

\paragraph{Inner instantiations.}
We likewise give three inner encoders, organized by work and by the strength of the contraction bound \(p_w(\delta)\).
\begin{enumerate}
\item \textbf{Accumulator inner.}
The simplest recursive rate-$1$ inner.  It admits an exact input--output enumerator via a spacing/compositions count.
This provides a clean baseline proof vehicle, but its certified \(\delta\) is limited.
\item \textbf{Dense time-varying recursive inner (log memory).}
A log-memory streaming recursion with fresh dense random taps at each time step.
While the encoder is \(O(n\log n)\)-work, it yields a strong three-mechanism tail bound for \(p_w(\delta)\)
(early-first-one, rare ON\(\to\)OFF termination, and a large-deviation term),
and a forced-termination refinement that turns termination into a geometric-in-\(w\) contraction for small/intermediate weights.
\item \textbf{Sparse-tap recursive inner with periodic refresh (linear work).}
A streaming recursion that uses only \(d=O(1)\) taps per step and applies a periodic linear refresh on the state.
We prove a refresh-to-balance property, exchangeable period starts, and between-refresh stability, giving a linear-work
inner that supports the same modular framework.
\end{enumerate}

\paragraph{Putting the modules together.}
The framework immediately yields a family of end-to-end concatenated ensembles, obtained by pairing any outer instantiation
with any inner instantiation.  The most practically relevant combination is the linear-work pairing:
banded-LDPC outer + sparse-refresh inner, which achieves linear encoding work with a single random interleaver
and provable linear minimum distance.

\paragraph{Organization.}
Section~\ref{sec:sc-framework} states the modular serial-concatenation framework theorem and isolates the required outer and inner interfaces.
Sections~\ref{sec:outer-dense-logmem} and \ref{sec:outer-local-expander} instantiate the outer interface (dense log-memory and banded LDPC).
Sections~\ref{sec:inner-accumulator}, \ref{sec:inner-dense}, and \ref{sec:sparse-refresh-inner} instantiate the inner interface
(accumulator, dense time-varying recursion, and sparse recursion with periodic refresh).
Section~\ref{sec:integrations} begins the end-to-end synthesis, starting with a dense outer + dense inner validation theorem
and setting up the later flagship sparse+sparse integration.
